% Exact staging-only source-aware proof for production Rethlas verification.
% Date: 29 August 2026.
% This file is not a live-book input and makes no novelty or priority claim.

\begin{proof}[Proof of
\cref{prop:p062-proper-full-rank-moment}]
We first verify that every current and test form used below is defined.
The function $h=\varphi\circ F$ is continuous.  If $\omega$ is a smooth
compactly supported $(n,n)$-form on $X$, then $h$ is bounded on
$\operatorname{Supp}(\omega)$, while $\omega$ defines a finite signed
Radon measure by \cite[\S\S5.8.5--5.8.6]{CLD25}.  Thus $h$ is locally
integrable and the integration current $[h]$ is defined by
\cite[Corollary~7.2.10]{CLD25}.

The sign convention requires one explicit translation.  Chambert--Loir and
Ducros use the tropical coordinates
$\tau_i=\log|f_i|$, whereas the statement uses $F_i=-\log|f_i|$.
Equivalently, $F$ is the tropical chart of the inverse-coordinate moment
$(f_1^{-1},\ldots,f_n^{-1})$; the affine isomorphism $L$ preserves
pleasantness.
For a Lagerberg form $\alpha$ in the $F$-coordinates, write
\begin{equation}
\label{eq:p062-moment-signed-pullback}
                         F^*\alpha:=\tau^*(L^*\alpha).
\end{equation}
The affine pullback $L^*$ commutes with $\mathrm d'$ and $\mathrm d''$,
and pullback by a tropical chart does as well, by
\cite[\S\S3.3.11--3.3.12]{CLD25}.  Moreover, the calibration is a positive
vector-volume, so
\begin{equation}
\label{eq:p062-moment-sign-calibration}
                           L_*e=e.
\end{equation}
There is no factor $(-1)^n$: equivalently, the sign $-1$ occurs once in
the primed and once in the double-primed top-degree factors.

For every compactly supported smooth Lagerberg form $\gamma$ on $\Omega$,
of any bidegree, the form $F^*\gamma$ in
\eqref{eq:p062-moment-signed-pullback} is an ordinary CLD-smooth form.  Its
support satisfies
\begin{equation}
\label{eq:p062-moment-test-support}
 \operatorname{Supp}(F^*\gamma)
 \subseteq F^{-1}\bigl(\operatorname{Supp}(\gamma)\bigr).
\end{equation}
The right-hand side is compact by
\cref{itm:p062-moment:proper}, since
$\operatorname{Supp}(\gamma)$ is compact in $\Omega$.  Hence every
$F^*\gamma$ is a legitimate compactly supported test form; no cutoff is
inserted.  The tropical chart and the affine map $L$ make $F$ a morphism of
the underlying tropical spaces, and
\cref{itm:p062-moment:proper} is precisely the topological properness used
in \cite[\S7.1.11]{CLD25}.  Thus the pushforward of currents $F_*$ is
defined.

The continuous function $\varphi$ is locally Lebesgue integrable on the
Euclidean open set $\Omega$, so the target current $[\varphi]$ is defined.
Let $\gamma\in A_c^{n,n}(\Omega)$ be arbitrary.  The top-degree form
$\varphi\gamma$ has continuous, hence Borel, coefficients and compact
support, and is integrable.  Here
$F^*(\varphi\gamma):=\tau^*L^*(\varphi\gamma)$ is the Borel-coefficient
pullback of \cite[Proposition~5.10.1]{CLD25}; it is not asserted to be
smooth.  The explicit calibration
\eqref{eq:p062-moment-explicit-pushforward-calibration}, the full-image
identity \eqref{eq:p062-moment-full-image}, and
\cite[Proposition~5.10.1(c),(d)]{CLD25} give
\begin{align}
 \bigl\langle F_*[h],\gamma\bigr\rangle
  &=\int_X(\varphi\circ F)F^*\gamma
    =\int_XF^*(\varphi\gamma)                         \notag\\
  &=\int_\Omega\langle\varphi\gamma,\mu_F\rangle
    =m\int_\Omega\varphi\gamma
    =\bigl\langle m[\varphi],\gamma\bigr\rangle.
\label{eq:p062-moment-pushforward-current-pairing}
\end{align}
Since $\gamma$ was arbitrary,
\begin{equation}
\label{eq:p062-moment-pushforward-current}
                              F_*[h]=m[\varphi].
\end{equation}

Pushforward by a proper morphism commutes with $\mathrm d'$ and
$\mathrm d''$ by \cite[\S7.7.8]{CLD25}.  Hence
\begin{equation}
\label{eq:p062-moment-target-zero-pushforward}
 m\,\mathrm d'\mathrm d''[\varphi]
 =\mathrm d'\mathrm d''F_*[h]
 =F_*\bigl(\mathrm d'\mathrm d''[h]\bigr)=0.
\end{equation}
Since $m>0$, this proves
\begin{equation}
\label{eq:p062-moment-target-zero-current}
                      \mathrm d'\mathrm d''[\varphi]=0
                      \quad\text{on }\Omega.
\end{equation}

It remains to use the Euclidean recognition theorem at exactly its stated
strength.  Let $B$ be an open Euclidean ball with compact closure in
$\Omega$.  The distribution-coefficient description of Euclidean currents
in \cite[\S1.6.4]{CLD25} identifies
$\mathrm d'\mathrm d''[\varphi]$ with the distributional Hessian of
$\varphi$ (up to the fixed coefficient signs); thus
\eqref{eq:p062-moment-target-zero-current} is the vanishing of every second
distributional derivative.  Equivalently, on $B$, both
$\mathrm d'\mathrm d''[\varphi]$ and
$\mathrm d'\mathrm d''[-\varphi]$ are the zero current and hence are
positive.  The equivalence of current positivity and convexity on a convex
Euclidean open set in \cite[Proposition~8.3.2]{CLD25} shows that
$\varphi|_B$ and $-\varphi|_B$ are convex.  Therefore $\varphi|_B$ is
affine.  Thus $\varphi$ is locally affine on $\Omega$.  Its locally defined
gradient is a locally constant map from the connected set $\Omega$ to
$\mathbf R^n$, hence is a single vector $(a_1,\ldots,a_n)$.  The function
$\varphi-\sum_i a_i y_i$ is locally constant and therefore constant on
$\Omega$.  This proves \eqref{eq:p062-moment-affine-target}, and pullback
gives \eqref{eq:p062-moment-affine-source}.

Each $\log|f_i|$ is an ordinary CLD-smooth function, and
\eqref{eq:p062-moment-affine-source} has zero
$\mathrm d'\mathrm d''$ as an ordinary smooth form.  By
\cite[Definition~7.2 and Remark~7.3(5),(6)]{GJR26}, logarithms of analytic
units are GJR harmonic and their real linear combinations are harmonic.
The constant term is covered as well: nontriviality of the valuation gives
$c\in K^\times$ with $\log|c|\ne0$, so
$b=(b/\log|c|)\log|c|$.  Thus $h$ is GJR harmonic and, by the same
definition, $\mathbf R$-piecewise linear.  Finally,
\cite[Proposition~7.2(2),(3)]{GR25Classical} says that logarithms of analytic
units are classically pluriharmonic and that classically pluriharmonic
functions form a real vector space.  The preceding scalar-unit observation
also covers the arbitrary constant $b$.  Hence $h$ is classically
pluriharmonic.

Suppose now that $X$ is boundaryless and that
\eqref{eq:p062-moment-affine-target} holds.  The source formula
\eqref{eq:p062-moment-affine-source} follows.  The Poincar\'e--Lelong formula
\cite[Theorem~16.7.3]{CLD25} applied to each analytic unit $f_i$ gives
\begin{equation}
\label{eq:p062-moment-unit-zero-current}
                 \mathrm d'\mathrm d''[\log|f_i|]
                 =\delta_{\operatorname{div}(f_i)}=0.
\end{equation}
Moreover, $[b]=b\delta_X$ and $\delta_X$ is
$\mathrm d'$- and $\mathrm d''$-closed by \cite[\S7.7.4]{CLD25}.
Linearity in \eqref{eq:p062-moment-affine-source} therefore gives
\eqref{eq:p062-moment-zero-current}.  This proves the asserted converse.
\end{proof}

\begin{proof}[Proof of
\cref{cor:p062-proper-full-rank-moment-automatic-calibration}]
The retained pleasantness hypothesis is exactly
\cite[Definition~14.1.7]{CLD25}; the affine isomorphism $L$ makes $F$ a
pleasant tropical chart as in \cite[Definition~5.4.3]{CLD25}.  Since $X$
is boundaryless, the harmony theorem \cite[Theorem~6.6.8]{CLD25} says
directly that $(\Pi_F,\mu_F)=(\Omega,\mu_F)$ is harmonious everywhere.

The calibration $\mu_{\Sigma_F^{(n)}}$ is standard by
\cite[Definition~14.3.5]{CLD25}, and
\cite[Definition~14.5.1]{CLD25} defines its image calibration $\mu_F$ and
states that it is standard.  By \cite[Definition~4.5.1]{CLD25}, a standard
calibration is strictly positive.  Choose a locally finite face-to-face
cellular decomposition of the open polyhedral set $\Omega$, adapted to
that calibration.  On every $n$-cell $C$ there is a number $m_C>0$ such
that
\begin{equation}
\label{eq:p062-moment-cell-weight}
                              \mu_F|_C=m_Ce.
\end{equation}
Positivity and nonvanishing of $m_C$ also follow directly from the local
degree description in \cite[\S14.5.2]{CLD25}.

Let $D$ be a codimension-one cell and take a point in its relative interior.
Because $\Omega$ is open in $\mathbf R^n$, exactly two local
$n$-dimensional sides, say $C_+$ and $C_-$, meet there.  Their primitive
outward normal directions are opposite.  The vanishing of the discordance
at $D$, which is the harmony condition, reads
\begin{equation}
\label{eq:p062-moment-facet-balance}
                         (m_{C_+}-m_{C_-})\nu_D=0.
\end{equation}
Hence $m_{C_+}=m_{C_-}$.
This is precisely the two-cell calculation isolated in
\cite[Remark~4.6.10]{CLD25}.

Any two top-dimensional cells can be connected by a polygonal path in the
connected open set $\Omega$.  After an arbitrarily small perturbation on
the compact image of the path, local finiteness lets us arrange that the
path avoids the codimension-two skeleton, crosses codimension-one cells
transversely, and crosses only finitely many of them.  Repeated application
of \eqref{eq:p062-moment-facet-balance} shows that all the numbers $m_C$
are one constant $m>0$.  Equation
\eqref{eq:p062-moment-cell-weight} therefore gives
\eqref{eq:p062-moment-constant-calibration}.  The corollary now follows
from \cref{prop:p062-proper-full-rank-moment}.
\end{proof}

\begin{remark}
\label{rem:p062-proper-full-rank-moment-boundary-audit}
The proof of the proposition uses no Stokes integration by parts and never
identifies $\mathrm d'\mathrm d''[h]$ with the current associated to a
smooth form.  It uses the definition of proper current pushforward in
\cite[\S7.1.11]{CLD25} and its differential functoriality in
\cite[\S7.7.8]{CLD25}; compact support in the open target removes target
boundary terms, and properness supplies compact support upstairs.  In
particular, boundarylessness is not needed for the forward implication in
the principal proposition.  That hypothesis is load-bearing both for its
Poincar\'e--Lelong converse and in the corollary,
where \cite[Theorem~6.6.8]{CLD25} would otherwise allow a discordance over
the image of the source boundary.
\end{remark}
